% TODO - formatar esse conteúdo aqui
\section{Probabilidade}
\label{cap:probabilidade}

Para o primeiro experimento, foi considerado o jogo Cara ou Coroa. Nesse jogo, um jogador
escolhe "Cara" e o outro "Coroa", em seguida uma moeda é lançada ao alto e é verificada
qual face cai virada para cima. Nesse caso, a chance de vir "Cara" é:
\[ P(Cara)=\frac{1}{2}=0,5 \]
Ao lançar 3 moedas, a chance de vir cara em todas é:
\[ P(Cara,Cara,Cara)=\frac{1}{2}\cdot\frac{1}{2}\cdot\frac{1}{2}=0,125 \]
pois são eventos independentes.

O código realiza o experimento $n$ vezes e calcula a probabilidade experimental (o total de
vezes que o evento ocorre dividido pelo total de experimentos). Após vários experimentos,
são calculados a média e variância dos resultados. A seguir está a implementação do
experimento:

\begin{lstlisting}[language=Java]

import java.util.Random;

public class Probabilidade {
    public static void main(String[] args) {
        double media = 0, variancia = 0, soma_desvio = 0;
        int quant_experimentos = 10;
        double eventos_total = 100;
        Random gerador = new Random();

        // cria um vetor para guardar a probabilidade de cada experimento
        double[] probabilidade = new double[quant_experimentos];

        for (int i = 0; i < quant_experimentos; i++) {
            // realiza o experimento i
            int eventos_sucesso = 0;
            for (int j = 0; j < eventos_total; j++) {
                // gerar 3 valores aleatórios de 0 ou 1
                // 0 = cara, 1 = coroa
                int r1 = gerador.nextInt(2);
                int r2 = gerador.nextInt(2);
                int r3 = gerador.nextInt(2);

                // verifica se as 3 moedas deram cara
                if (r1 == 0 && r2 == 0 && r3 == 0) {
                    eventos_sucesso++;
                }
            }
            // calcula a probabilidade experimental e guarda esse valor
            probabilidade[i] = eventos_sucesso / eventos_total;
            // soma para fazer a média
            media += probabilidade[i];
        }

        // calcula a média
        media = media / quant_experimentos;

        // calcula a variância
        for (int i = 0; i < quant_experimentos; i++) {
            double desvio = probabilidade[i] - media;
            desvio = desvio * desvio;
            soma_desvio += desvio;
        }

        // calcula a variância usando os desvios
        variancia = soma_desvio / quant_experimentos;

        // indica as métricas
        System.out.println("Métricas:");
        System.out.println("Para " + quant_experimentos + " experimentos, com " + eventos_total + " eventos cada");
        System.out.println("Média: " + media);
        System.out.println("Variância: " + variancia);
    }
}

\end{lstlisting}


